\documentclass[10pt,a4paper]{amsart}
\usepackage[margin=19mm]{geometry}
\usepackage{amsmath,amssymb,mathtools}
\usepackage[scr=rsfs,cal=euler]{mathalfa}
\usepackage{xfrac}
\usepackage[unicode,hidelinks,bookmarks=false]{hyperref}
\newcommand{\ie}{i.e.\ }
\newcommand{\cf}{cf.\ }
\newcommand{\resp}{respectively\ }
\newcommand{\Subsection}{\subsection}
\providecommand{\nobreakdash}{\nobreak\hbox{-}}
\setlength{\emergencystretch}{2em}
\multlinegap0pt
\usepackage{bm}
\usepackage{bbm}
\usepackage{dsfont}
\usepackage{xspace}
\usepackage{cancel}
\usepackage{mathtools}
\usepackage{amsthm}
\newtheorem{theorem}{Theorem}
\newtheorem{lemma}[theorem]{Lemma}
\newcommand{\ind}{\mathbf 1}
\title[Extremal expectations in card guessing with partial feedback]{Extremal expectations in card guessing with partial feedback: proofs of two conjectures of Diaconis-Graham-Spiro}
\author{Congyi Luo}
\date{Source: arXiv:2609.25062v1 (CC BY 4.0)}
\begin{document}
\maketitle

\noindent\emph{Source and licence.} Extremal expectations in card guessing with partial feedback: proofs of two conjectures of Diaconis-Graham-Spiro. Version arXiv:2609.25062v1 (CC BY 4.0), distributed under the Creative Commons Attribution 4.0 International licence, as recorded in the arXiv licence metadata for this exact version and shown on the abstract page https://arxiv.org/abs/2609.25062v1. Full rights notice: licenses/arxiv-2609.25062-CC-BY-4.0.txt.\par\smallskip

\noindent\textit{Editorial scope.} Counted unit: the Lemma whose source label is \texttt{low:lem:switch} in the article (source0.tex lines 248--253, source label \texttt{low:lem:switch}), delivered together with the article's own complete original proof of it (source0.tex lines 254--284, one \texttt{proof} environment of 31 lines). No mathematics is written, repaired or re-derived: statement and proof are the article's own text line for line. Required context retained uncounted: none. Every definition, display and label of those lines is retained unchanged. Omitted from this package and not redistributed: 1--247, 285--635. Nothing inside a retained excerpt is omitted or altered: every retained body is delivered exactly as the article's lines are written, so no omission receipt of any kind accompanies this package. Citations: the retained excerpts carry no citation command at all, so no bibliography entry of the article is transcribed and no reference is redistributed. Reference adaptation: none was needed and none was made; every reference inside the delivered unit resolves inside it, so no unresolved reference or citation is produced. Printed slips kept literal: no printed word, symbol, hypothesis or number of the counted statement or of its proof was altered. Layout and class: the article's own class is \texttt{amsart} with packages including geometry, amsmath, amssymb, amsthm, booktabs, hyperref; the wrapper is the lane's pinned \texttt{amsart} editorial wrapper at 10pt on A4, which restates the theorem-like environments the retained text needs and adds the article's own macro definitions that the retained text needs (\texttt{\textbackslash ind}), with extra wrapper packages (bm, bbm, dsfont, xspace, cancel, mathtools). The delivered numbering is the unit's own. Assets: no retained span contains a figure environment, a graphics-inclusion command, a PDF-page inclusion, a table or a TikZ picture; no image file is copied into this package, the package declares no asset dependency and no image file is redistributed with it. Licence: version arXiv:2609.25062v1 of this article is distributed under the Creative Commons Attribution 4.0 International licence, as recorded in the arXiv licence metadata for this exact version and shown on the abstract page https://arxiv.org/abs/2609.25062v1. A full rights notice is delivered at licenses/arxiv-2609.25062-CC-BY-4.0.txt. Characters: every retained excerpt is pure ASCII, so the delivered engine, XeLaTeX with its default Latin Modern text font, reports no missing character.
\begin{lemma}\label{low:lem:switch}
If $L\ge m$, then for every $i\in[n]$,
\[
 p_i\ge\left(1-\frac mL\right)\frac{r_i}{M-f_i}.
\]
\end{lemma}

\begin{proof}
Every free position allows $i$, so $M-f_i\ge L>0$. If $L=m$, the right-hand side is zero and the result follows. Assume henceforth that $L>m$, and fix a free position $z$.

Let $x$ be a position forbidding $j$, where $j\ne i$. In $\mathcal W$, consider
\[
 A_x=\{w_x=i,\ w_z\ne j\},\qquad B_z=\{w_z=i\}.
\]
Swapping positions $x,z$ is a bijection from $A_x$ to $B_z$. Indeed, for $w\in A_x$, the new label at $x$ is not $j$, and the new label at $z$ is $i$. All other positions and all multiplicities are unchanged, so the resulting word belongs to $\mathcal W$ and to $B_z$. Conversely, for $w\in B_z$, its original label $w_x$ is not $j$. After the swap, $x$ contains the allowed label $i$ and $z$ contains the original $w_x\ne j$, so the word belongs to $A_x$. Applying the swap twice restores the word. Uniformity gives
\[
 p_i=\mathbb P(A_x).
\]
If $\mathbb P(w_x=i)=0$, the inequality $p_i\ge(1-m/L)\mathbb P(w_x=i)$ is immediate. Otherwise, conditional on $w_x=i$, the free positions are still symmetric, and at most $r_j$ of them contain $j$. Hence
\[
 L\mathbb P(w_z=j\mid w_x=i)
 =\mathbb E\left[\sum_{y\text{ free}}\ind\{w_y=j\}\,\middle|\,w_x=i\right]
 \le r_j\le m.
\]
Thus
\[
 p_i=\mathbb P(w_x=i)
       \bigl(1-\mathbb P(w_z=j\mid w_x=i)\bigr)
 \ge(1-m/L)\mathbb P(w_x=i).
\]
If $x$ is itself free, then $\mathbb P(w_x=i)=p_i$ and the same comparison holds. If $x$ forbids $i$, that probability is zero. Let $I_i$ be the set of positions allowing $i$. Summing the comparison over $x\in I_i$ gives
\[
 (M-f_i)p_i
 \ge(1-m/L)\sum_{x\in I_i}\mathbb P(w_x=i)
 =(1-m/L)r_i.
\]
The last equality holds because each word contains exactly $r_i$ copies of $i$, all in $I_i$. Division by $M-f_i>0$ proves the claim.
\end{proof}

\end{document}
